% =====================================================================
%  Compositional GMM-Router (CGR) -- Implementation Specification
%  Build: pdflatex compositional_gmm_router_spec.tex   (run twice)
% =====================================================================
\documentclass[11pt]{article}
\usepackage[margin=2.2cm]{geometry}
\usepackage{amsmath,amssymb,bm}
\usepackage{booktabs,tabularx,array}
\usepackage{algorithm}
\usepackage{algpseudocode}
\usepackage{enumitem}
\usepackage{xcolor}
\usepackage{listings}
\usepackage{hyperref}
\hypersetup{colorlinks=true,linkcolor=blue!50!black,urlcolor=blue!50!black,citecolor=blue!50!black}

% ---------- macros ----------
\newcommand{\R}{\mathbb{R}}
\newcommand{\T}{^{\top}}
\newcommand{\tagthesis}{\textsuperscript{\textcolor{teal!70!black}{\sffamily[thesis]}}}
\newcommand{\tagnew}{\textsuperscript{\textcolor{orange!80!black}{\sffamily[new]}}}
\newcommand{\tagassumed}{\textsuperscript{\textcolor{red!70!black}{\sffamily[assumed]}}}
\newcommand{\MUST}{\textbf{MUST}}
\newcommand{\MUSTNOT}{\textbf{MUST NOT}}
\newcommand{\SHOULD}{\textbf{SHOULD}}
\newcommand{\MAY}{\textbf{MAY}}
\newcommand{\Ladapt}{\mathcal{L}_{\mathrm{ad}}}
\newcommand{\role}{\par\noindent\textbf{Role.}\ }
\newcommand{\impl}{\par\noindent\textbf{Implementation requirements.}\ }

\lstset{
  language=Python,
  basicstyle=\ttfamily\footnotesize,
  keywordstyle=\color{blue!60!black}\bfseries,
  commentstyle=\color{green!40!black},
  stringstyle=\color{red!50!black},
  frame=single,
  breaklines=true,
  columns=fullflexible,
  showstringspaces=false
}

\algnewcommand{\LineComment}[1]{\State \(\triangleright\) #1}

\title{\textbf{Compositional GMM-Router (CGR)}\\[4pt]
\large Leak-free routing with learned cross-task transfer for continual code generation\\[2pt]
\normalsize Implementation specification}
\author{}
\date{October 2026}

\begin{document}
\maketitle

\begin{abstract}
\noindent This document fully specifies the Compositional GMM-Router (CGR), an extension of the GMM-Router of the
graduation thesis \emph{Tran An Khanh (HUST, 2026)}. CGR removes the thesis's iterative refinement pipeline
(calibration inference, traceback router, adapter refinement on repair pairs, JSD-gated Round-2), because those
stages pass ground-truth task identity to later stages through the compiler and the test function. CGR adds
explicit, learned knowledge transfer between tasks while keeping per-task isolation of trained parameters. The
document is written so that an implementing agent can build, test and evaluate the method from this file alone.
\end{abstract}

\tableofcontents
\newpage

% =====================================================================
\section{How to use this document}
% =====================================================================

\paragraph{Audience.} An implementing agent (e.g.\ Claude Code) or engineer. Every component is given as
(i)~a formulation, (ii)~its role in the method, and (iii)~implementation requirements.

\paragraph{Normative language.} \MUST\ and \MUSTNOT\ are hard requirements; violating them invalidates the
method (in particular the leak-free property). \SHOULD\ is a strong default that may be changed only with a
logged reason. \MAY\ marks optional optimisations.

\paragraph{Provenance tags.} Values tagged \tagthesis\ are taken from the GMM-Router thesis and keep the setup
comparable. Values tagged \tagnew\ are introduced by CGR and are tunable. Values tagged \tagassumed\ are not
stated in the thesis; the implementer \MUST\ check them against the thesis codebase if it is available and
record the value actually used.

\paragraph{Relationship to the thesis codebase.} If the GMM-Router codebase is available, the implementer
\SHOULD\ reuse its data loading, instruction templates, tokenisation, decoding settings and evaluation harness
(BLEU, CodeBLEU, Pass@$k$ execution), and \MUST\ replace its router and adapter modules with the components of
Section~\ref{sec:components}. The following thesis code paths \MUST\ be deleted or made unreachable from CGR:
calibration inference (thesis Step~2), adapter refinement (Step~3a), traceback GMM fitting (Step~3b), Round-2
inference and JSD-gated score combination (Step~5). The calibration split is not used by CGR.

\paragraph{Build order (milestones).}
\begin{enumerate}[label=M\arabic*., leftmargin=*]
  \item Routing encoder, diagonal EM, within-task scatter and router scoring (Sections~\ref{sec:embed}--\ref{sec:score}).
        Pass tests T1--T4 (Section~\ref{sec:tests}).
  \item Atom bank and composition matrix with $C=I$; confirm that CGR then reduces to thesis GMM-Router Round-1
        (merge equivalence, test T5).
  \item Transfer prior, $\beta$ initialisation and two-phase training (Sections~\ref{sec:prior}--\ref{sec:train});
        pass tests T6--T7.
  \item Inference API with leakage and no-replay guards (Section~\ref{sec:infer}); pass tests T8--T10.
  \item Full task sequences on both benchmarks, evaluation matrix, continual-learning metrics, and the ablation
        grid of Section~\ref{sec:ablations}.
\end{enumerate}

% =====================================================================
\section{Method summary}
% =====================================================================

\noindent\textbf{Title:} Compositional GMM-Router

\medskip
\noindent\textbf{Research hypothesis:} Separating task identification from knowledge reuse can give cross-task
transfer comparable to parameter-sharing methods while keeping the near-zero forgetting of per-task isolation.
Identification means routing each query from its input alone with a shared-covariance GMM. Reuse means building
each new expert as a learned combination of frozen earlier LoRAs plus a fresh residual. Neither part needs
previous-task data, execution feedback, or extra forward passes at inference.

\medskip
\noindent For each arriving task, we use the existing GMM router to measure how strongly the new task's inputs
resemble each earlier task. We initialize one lateral weight per earlier LoRA from that affinity, then train these
weights together with a fresh residual LoRA on the new task's data while all earlier LoRAs stay frozen. At
inference, routing weights computed from the input alone pass through the learned lateral weights to merge every
LoRA into one adapter. The output is generated in a single round with no calibration, traceback, compiler or test
function.

\paragraph{Pipeline at a glance.} For each task $k$: (1)~embed the task's prompts with the frozen backbone;
(2)~compute the router's affinity of task $k$ to earlier tasks and initialise lateral weights; (3)~fit and
register task $k$'s GMM and add its within-task scatter to the shared covariance; (4)~calibrate the routing
temperature on a held-out slice of task $k$; (5)~Phase~A: train lateral weights only; (6)~Phase~B: train a new
LoRA atom and the lateral weights; (7)~freeze everything belonging to task $k$. At inference: embed
$\rightarrow$ score $\rightarrow$ route $\rightarrow$ compose $\rightarrow$ generate once.

% =====================================================================
\section{Hard constraints (invariants)}\label{sec:invariants}
% =====================================================================

\begin{enumerate}[label=I\arabic*., leftmargin=*]
  \item \textbf{No task identity at inference.} The prediction function \MUST\ accept only prompt strings.
        It \MUSTNOT\ receive or look up task index, language label, dataset name, test cases or any other metadata.
        Task labels may be attached to outputs only \emph{after} generation, for evaluation.
  \item \textbf{Single generation round.} Each query is generated exactly once (with $n$ samples drawn in that one
        round for Pass@$k$). There is no iterative refinement, no Round-2, no regeneration conditioned on earlier
        outputs, and no re-ranking or selection of candidates by any execution signal.
  \item \textbf{No execution feedback anywhere in the method.} No compiler, interpreter, test harness, runtime
        error message or traceback is used in training, routing, calibration or inference. Execution happens only
        in the evaluation harness, which reads generation files after all generation has finished and writes
        nothing back into the method's state.
  \item \textbf{No previous-task data.} After task $k$ is finished, its raw data, prompts, targets and routing
        embeddings \MUSTNOT\ be used again by the method. Allowed persistent state is exactly the list in
        Section~\ref{sec:state}. (Test sets of earlier tasks are read only by the evaluation harness.)
  \item \textbf{Frozen backbone and frozen past.} Backbone weights are never updated. After task $k$ is finished,
        atom $k$, row $k$ of every composition matrix, and GMM $k$ are frozen forever.
  \item \textbf{Adapter-independent routing features.} Routing embeddings are computed from blocks that carry no
        adapters (or, equivalently, with all adapters disabled).
  \item \textbf{Inputs as given.} Prompts are used exactly as produced by the benchmark templates. No extra prefix
        or token encoding task identity may be added beyond what the benchmark's own instruction contains.
\end{enumerate}

% =====================================================================
\section{Notation and persistent state}\label{sec:state}
% =====================================================================

Tasks $T_1,\dots,T_K$ arrive in order; $D_k=\{(x_i,y_i)\}$ is the training set of task $k$ (prompt $x$, target $y$).
$f_\theta$ is the frozen backbone with $L$ decoder blocks indexed $0,\dots,L-1$ and hidden size $d$.
$\ell^*$ is the routing/split block. The adapted blocks are $\Ladapt=\{\ell^*,\dots,L-1\}$ and the adapted modules
in each block are $u\in\{\texttt{q\_proj},\texttt{v\_proj}\}$. $p$ is the routing dimension, $M$ the number of
GMM components per task, $r$ the LoRA rank.

\begin{table}[h]
\centering
\small
\begin{tabularx}{\textwidth}{@{}l l X@{}}
\toprule
Symbol & Shape & Meaning (all items are frozen once written, except the running scatter and $\tau$)\\
\midrule
$P$ & $p\times d$ & Fixed orthonormal routing projection, $PP\T=I_p$.\\
$(\pi_{km},\mu_{km},\sigma^2_{km})$ & $M$, $M\times p$, $M\times p$ & Diagonal GMM of task $k$ (weights, means, variances).\\
$S_k$, $N_k$ & $p\times p$, scalar & Within-task scatter of task $k$ and number of embeddings used.\\
$S$, $n$ & $p\times p$, scalar & Running sums $S=\sum_k S_k$, $n=\sum_k N_k$; $\bar\Sigma=S/n$.\\
$(A_k^{(\ell,u)},B_k^{(\ell,u)})$ & $r\times d_{\mathrm{in}}$, $d_{\mathrm{out}}\times r$ & LoRA \emph{atom} of task $k$ at block $\ell\in\Ladapt$, module $u$.\\
$\beta_{k}^{(\ell)}$ & $k-1$ per block & Lateral weights of task $k$ on earlier atoms at block $\ell$.\\
$C^{(\ell)}$ & $K\times K$ & Lower-triangular composition matrix: $C_{kk}=1$, $C_{kj}=\beta^{(\ell)}_{kj}$ for $j<k$, $0$ above the diagonal.\\
$\tau$ & scalar & Routing temperature (running calibrated value, Section~\ref{sec:score}).\\
\bottomrule
\end{tabularx}
\caption{Persistent state of CGR. Nothing else may be stored across tasks.}
\end{table}

% =====================================================================
\section{Components}\label{sec:components}
% =====================================================================

\subsection{Backbone, adapted layers and routing embedding}\label{sec:embed}

\paragraph{Formulation.} For a prompt $x$ with $T$ tokens, let $H\in\R^{T\times d}$ be the hidden states after
blocks $0,\dots,\ell^*-1$ (Hugging Face convention: \texttt{output\_hidden\_states[}$\ell^*$\texttt{]}, where index~0
is the embedding output)\tagassumed. With token mask $m_t\in\{0,1\}$ (1 for non-padding prompt tokens):
\begin{align}
  \tilde h &= \frac{\sum_{t=1}^{T} m_t H_t}{\sum_{t=1}^{T} m_t}\in\R^{d}, \\
  h &= \frac{\tilde h-\mathrm{mean}(\tilde h)\,\mathbf 1}{\mathrm{std}(\tilde h)+\epsilon_{\mathrm{LN}}}\in\R^{d}
  \qquad\text{(parameter-free layer norm over the $d$ entries)}, \\
  z(x) &= P\,h\in\R^{p}, \qquad P=Q\T,\quad QR=\mathrm{qr}(G),\quad G\in\R^{d\times p},\ G_{ij}\sim\mathcal N(0,1)\ \text{(seed $s_P$)}.
\end{align}

\role A fixed feature space that sees only the query. It is computed before any token is generated, so nothing
generated, compiled or executed can influence routing (Invariant~I6). Because $P$ never changes, every task's
GMM stays valid as tasks accumulate.

\impl
\begin{itemize}[leftmargin=*]
  \item \MUST\ use the same prompt string and template as used for generation; prompt tokens only (no target
        tokens); right truncation at 2048 tokens\tagthesis.
  \item \MUST\ compute $h$ and $z$ in float32; $\mathrm{std}$ is the population standard deviation; $\epsilon_{\mathrm{LN}}=10^{-8}$.
  \item \MUST\ generate $P$ once from seed $s_P$, save it, and reload it; it is never regenerated.
  \item LoRA atoms are installed only on blocks $\ell\ge\ell^*$, so blocks $<\ell^*$ are adapter-free. The
        implementer \MAY\ therefore read $H$ from the generation prefill instead of a separate forward pass, but test
        T1 \MUST\ show that the result equals the adapters-disabled computation.
\end{itemize}

\subsection{Per-task diagonal GMM and within-task scatter}\label{sec:gmm}

\paragraph{Formulation.} When task $k$ arrives, let $Z_k=\{z_i\}_{i=1}^{N_k}$ be routing embeddings of up to
$N_{\max}=5000$\tagthesis\ prompts sampled (with a fixed seed) from the training part of $D_k$. Fit a diagonal
GMM $p_k(z)=\sum_{m=1}^{M}\pi_{km}\,\mathcal N\!\big(z;\mu_{km},\mathrm{diag}(\sigma^2_{km})\big)$ by EM exactly as in
the thesis (Algorithm~\ref{alg:em}): responsibilities
\begin{equation}
  r_{im}=\frac{\pi_{km}\,\mathcal N(z_i;\mu_{km},\mathrm{diag}\,\sigma^2_{km})}{\sum_{m'}\pi_{km'}\,\mathcal N(z_i;\mu_{km'},\mathrm{diag}\,\sigma^2_{km'})},
\end{equation}
M-step $\pi_{km}=N_m/N_k$, $\mu_{km}=\frac{1}{N_m}\sum_i r_{im}z_i$,
$\sigma^2_{kmj}=\max\!\big(\frac1{N_m}\sum_i r_{im}(z_{ij}-\mu_{kmj})^2,\ \sigma^2_{\min}\big)$ with $N_m=\sum_i r_{im}$.
After convergence, with responsibilities from a final E-step, compute the within-task scatter and update the
running sums:
\begin{equation}
  S_k=\sum_{i=1}^{N_k}\sum_{m=1}^{M} r_{im}\,(z_i-\mu_{km})(z_i-\mu_{km})\T,\qquad
  S\leftarrow S+S_k,\quad n\leftarrow n+N_k,\quad \bar\Sigma=S/n .
\end{equation}

\role GMM $k$ models where task $k$'s queries live in routing space. $S_k$ contributes task $k$'s within-task
spread to a single covariance $\bar\Sigma$ that is shared by all tasks; this is what removes the advantage that
broad task densities have over narrow ones (Section~\ref{sec:score}).

\impl
\begin{itemize}[leftmargin=*]
  \item Initialisation\tagthesis: $\pi_{km}=1/M$; $\mu_{km}$ sampled uniformly without replacement from $Z_k$;
        $\sigma^2_{km}$ set to the global per-dimension variance of $Z_k$ (floored at $\sigma^2_{\min}$).
  \item E-step \MUST\ use log-sum-exp; convergence when the change of mean log-likelihood is below
        $\delta=10^{-4}$ or after $T_{\max}=50$ iterations\tagthesis.
  \item $S_k$ \MUST\ be accumulated in float64. Store $S_k$ per task as well as the running $S$ (needed by
        ablations; it is a second-moment summary, not data).
  \item After this step, $Z_k$ \MUST\ be discarded (Invariant~I4).
\end{itemize}

\subsection{Shared-covariance scoring and routing weights}\label{sec:score}

\paragraph{Formulation.} With $K$ tasks registered and shrinkage $\lambda\in[0,1]$:
\begin{align}
  \Sigma_{km}(\lambda) &= (1-\lambda)\,\mathrm{diag}(\sigma^2_{km})+\lambda\big(\bar\Sigma+\epsilon_\Sigma I_p\big),\\
  \log\mathcal N(z;\mu,\Sigma) &= -\tfrac12\Big(p\log 2\pi+\log\det\Sigma+(z-\mu)\T\Sigma^{-1}(z-\mu)\Big),\\
  s_k(z) &= \operatorname{logsumexp}_{m}\Big[\log\pi_{km}+\log\mathcal N\big(z;\mu_{km},\Sigma_{km}(\lambda)\big)\Big]+\log\tfrac1K,\label{eq:score}\\
  \alpha_k(x) &= \frac{\exp\big(s_k(z(x))/\tau\big)}{\sum_{j=1}^{K}\exp\big(s_j(z(x))/\tau\big)},\qquad
  \hat k(x)=\arg\max_k s_k(z(x)).
\end{align}
Hard routing uses $\alpha(x)=e_{\hat k(x)}$. Setting $\lambda=0$ reproduces the thesis router (Eq.~4.6 of the thesis)
exactly.

\paragraph{Temperature calibration.} Raw score gaps are tens of nats (the thesis reports a 74.1-nat gap for one
query), so $\tau=1$ behaves like hard routing. CGR sets $\tau$ without execution and without previous-task data:
for each task $k\ge2$, using a held-out slice $H_k$ of $D_k$ (5\%, at most 500 prompts\tagnew) and the router
\emph{after} registering task $k$,
\begin{equation}
  \tau_k=\arg\min_{\tau\in\mathcal T}\ -\frac{1}{|H_k|}\sum_{x\in H_k}\log\alpha_k(x;\tau),\qquad
  \mathcal T=\{2^{-2},2^{-1},\dots,2^{8}\},\qquad
  \tau=\exp\Big(\tfrac{1}{k-1}\textstyle\sum_{j=2}^{k}\log\tau_j\Big).
\end{equation}

\role Task identification from the query alone; this replaces the role the traceback router played in the thesis,
without its leak. With $\lambda=1$, all tasks are measured in one metric, so a query far from every task is
assigned to the nearest task mean in whitened space instead of to the broadest density. The thesis attributes the
executable-benchmark collapse (15.8\% hard-routing accuracy, most queries routed to C++) to broad coverage by one
GMM; the shared metric targets exactly that mechanism. The same scores feed the transfer prior
(Section~\ref{sec:prior}).

\impl
\begin{itemize}[leftmargin=*]
  \item \MUST\ compute scores in float64 using a Cholesky factor of each $\Sigma_{km}(\lambda)$
        ($\log\det\Sigma=2\sum\log\mathrm{diag}\,L$; Mahalanobis term via a triangular solve).
  \item Because $\bar\Sigma$ changes when a task is registered, all $K\cdot M$ Cholesky factors \MUST\ be
        recomputed after each registration (stored GMM parameters do not change).
  \item For $\lambda=0$, \MUST\ use the plain diagonal formula (no $\epsilon_\Sigma$), so that test T4 holds exactly.
  \item On the CodeTask benchmark, hard routing is the default\tagthesis; soft routing with calibrated $\tau$ is
        reported as an ablation. On the executable benchmark, soft routing with calibrated $\tau$ is the default.
\end{itemize}

\subsection{LoRA atom bank and composition matrix}\label{sec:atoms}

\paragraph{Formulation.} Each task contributes one rank-$r$ LoRA \emph{atom} at every adapted block and module.
For a frozen linear map $W^{(\ell,u)}$ and coefficient vector $c^{(\ell)}\in\R^{K}$:
\begin{equation}
  y=W^{(\ell,u)}x+\kappa\sum_{i=1}^{K}c^{(\ell)}_i\,B_i^{(\ell,u)}A_i^{(\ell,u)}\,\mathrm{drop}(x),\qquad
  \kappa=\frac{\texttt{lora\_alpha}}{r}=2. \label{eq:atombank}
\end{equation}
The effective adapter of task $k$ is its own atom plus a learned combination of earlier atoms:
\begin{equation}
  \Delta W_k^{(\ell,u)}=B_k^{(\ell,u)}A_k^{(\ell,u)}+\sum_{j<k}\beta^{(\ell)}_{kj}B_j^{(\ell,u)}A_j^{(\ell,u)}
  =\sum_{i=1}^{K}C^{(\ell)}_{ki}\,B_i^{(\ell,u)}A_i^{(\ell,u)},
\end{equation}
so task $k$ is executed by setting $c^{(\ell)}=C^{(\ell)\T}e_k$ (row $k$ of $C^{(\ell)}$).

\role The explicit transfer channel. Each new task can reuse any earlier atom with a learned weight per block,
including weight zero or negative weight. Earlier atoms and earlier rows of $C$ are frozen, so reuse can never
change how an earlier task behaves; parameter-level forgetting stays at zero, as in the thesis.

\impl
\begin{itemize}[leftmargin=*]
  \item Atoms are installed on \texttt{q\_proj} and \texttt{v\_proj} of blocks $\ell\in\Ladapt$ only;
        $r=16$, $\texttt{lora\_alpha}=32$, dropout $0.1$ on the atom input, no bias\tagthesis.
  \item $\beta^{(\ell)}_{k}$ is one scalar per (adapted block, earlier atom), shared by \texttt{q\_proj} and
        \texttt{v\_proj} in that block\tagnew; unconstrained real values.
  \item \SHOULD\ implement Eq.~\eqref{eq:atombank} without materialising merged weights, by concatenating
        the atoms and scaling the rank dimension:
        \begin{gather*}
          A_{\mathrm{cat}}=[A_1;\dots;A_K]\in\R^{Kr\times d_{\mathrm{in}}},\qquad
          B_{\mathrm{cat}}=[B_1,\dots,B_K]\in\R^{d_{\mathrm{out}}\times Kr},\\
          y=Wx+\kappa\,B_{\mathrm{cat}}\big((c\otimes\mathbf 1_r)\odot A_{\mathrm{cat}}\,\mathrm{drop}(x)\big).
        \end{gather*}
        This supports a different $c$ per example in a batch and is mathematically identical to weight merging.
  \item New atom initialisation: $A_k$ Kaiming-uniform (PEFT default), $B_k=0$.
  \item Atoms and $\beta$ are kept in float32; the backbone in bf16\tagnew.
\end{itemize}

\subsection{Router-derived transfer prior and $\beta$ initialisation}\label{sec:prior}

\paragraph{Formulation.} \emph{Before} registering task $k$ ($k\ge2$), score its embeddings against the $k-1$
existing tasks only:
\begin{equation}
  a_{kj}=\frac{1}{N_k}\sum_{i=1}^{N_k}\frac{\exp\big(s_j(z_i)/\tau_{\mathrm{aff}}\big)}{\sum_{j'<k}\exp\big(s_{j'}(z_i)/\tau_{\mathrm{aff}}\big)},\qquad j<k,\qquad \tau_{\mathrm{aff}}=1,
\end{equation}
and initialise the lateral weights as the affinity-weighted merge of earlier \emph{effective} adapters:
\begin{equation}
  \beta^{(\ell)}_{kj}\leftarrow\sum_{i<k}a_{ki}\,C^{(\ell)}_{ij}\quad (j<k),\qquad\text{equivalently}\qquad
  \Delta W_k^{(\ell,u)}\big|_{\text{init}}=\sum_{i<k}a_{ki}\,\Delta W_i^{(\ell,u)}\quad(\text{with }B_k=0).
\end{equation}

\role The router's data-free guess of which earlier tasks are related to task $k$. Training starts from roughly
the adapter the current system would already apply to task $k$'s inputs, so learning can only move away from it
if task $k$'s data says so.

\impl For $k=1$ there are no earlier atoms and no $\beta$. For $k=2$, $a_{21}=1$.

\subsection{Two-phase training}\label{sec:train}

\paragraph{Formulation.} The loss for task $k$ is the masked negative log-likelihood of the target under task
$k$'s effective adapter:
\begin{equation}
  \mathcal L_k=-\frac{1}{|\mathcal B|}\sum_{(x,y)\in\mathcal B}\sum_{t=1}^{|y|}\log p\big(y_t\mid x,y_{<t};\ c^{(\ell)}=C^{(\ell)\T}e_k\ \forall\ell\big),
\end{equation}
with prompt and padding tokens masked out of the loss\tagthesis.
\begin{itemize}[leftmargin=*]
  \item \textbf{Phase A (compose).} Trainable: $\{\beta^{(\ell)}_k\}_{\ell\in\Ladapt}$ only. Atom $k$ is held at
        $B_k=0$ (contributes nothing). $S_A=300$ steps, batch 32, AdamW, $\eta_A=5\times10^{-3}$, no weight decay\tagnew.
        Skipped for $k=1$.
  \item \textbf{Phase B (specialise).} Trainable: atom $k$ ($A_k,B_k$) and $\beta_k$. Atom: AdamW,
        $\eta=10^{-4}$, weight decay $0.01$, warmup ratio $0.01$, 3 epochs, batch 32\tagthesis.
        $\beta_k$: learning rate $\eta_\beta=2\times10^{-5}$, no weight decay\tagnew.
  \item After Phase B, freeze atom $k$ and row $k$ of every $C^{(\ell)}$ (Invariant~I5).
\end{itemize}

\role The affinity is only a guess. Phase~A corrects it with task $k$'s own data: it keeps atoms that help and
pushes misleading ones toward zero or below. A typical misleading case is two tasks with similar inputs but
incompatible outputs, such as the Ruby inputs of CodeSearchNet (summarisation) and RunBugRun (repair). Phase~A's
loss is also a direct measurement of forward transfer, i.e.\ how much of task $k$ existing knowledge already
solves. Phase~B lets the new atom learn only what earlier atoms cannot express. Because $B_k$ starts at zero,
Phase~B starts exactly from the Phase-A composition (test T6).

\impl Use the final weights of each phase (no early stopping)\tagnew; log the loss on $H_k$ at the end of
Phase~A and every 500 steps of Phase~B. Maximum 100{,}000 training examples per task\tagthesis.

\subsection{Inference: composed merging and generation}\label{sec:infer}

\paragraph{Formulation.} For a query $x$:
\begin{equation}
  z=z(x),\qquad \alpha=\alpha(x)\ \text{(soft) or } e_{\hat k(x)}\ \text{(hard)},\qquad
  c^{(\ell)}(x)=C^{(\ell)\T}\alpha(x)\quad\forall\ell\in\Ladapt,
\end{equation}
then generate once with Eq.~\eqref{eq:atombank} using $c^{(\ell)}(x)$. By linearity,
\begin{equation}
  \sum_{i=1}^{K}c^{(\ell)}_i(x)\,B_iA_i=\sum_{k=1}^{K}\alpha_k(x)\,\Delta W_k^{(\ell,u)},
\end{equation}
which is the thesis's weight-level merge (thesis Eq.~4.10) applied to effective adapters. With $C=I$ CGR reduces
to thesis GMM-Router Round-1.

\role Delivers transfer at the same cost as the thesis's merging: one coefficient computation per query and one
forward pass per decoding step.
\emph{Forward transfer} enters through the routed task's row of $C$.
\emph{Backward transfer} enters only through routing: $c_j(x)=\sum_k\alpha_k(x)C_{kj}$ collects weight from later
tasks that reuse atom $j$, so an old-task query with some posterior mass on a later related task receives that
task's added knowledge. Changing an old task's own composition would require its data, which Invariant~I4 forbids.

\impl
\begin{itemize}[leftmargin=*]
  \item The public API \MUST\ be \texttt{predict(prompts, n\_samples)} with no other per-query inputs (Invariant~I1).
  \item All $n$ samples of a query share the same $c(x)$; no candidate is dropped, re-ranked or regenerated (I2).
  \item Decoding\tagthesis: CodeTask --- temperature 0.2, top-$p$ 0.95, at most 512 new tokens;
        executable --- $n=5$ samples, at most 2048 new tokens, temperature 0.2 and top-$p$ 0.95\tagassumed.
  \item Log $\alpha(x)$ and $\hat k(x)$ for every query (needed for routing diagnostics).
\end{itemize}

% =====================================================================
\section{Algorithms}
% =====================================================================

\begin{algorithm}[h]
\caption{\textsc{FitTaskGMM}: diagonal EM and within-task scatter (thesis Algorithm~1 plus scatter)}\label{alg:em}
\begin{algorithmic}[1]
\Require embeddings $Z_k=\{z_i\}_{i=1}^{N_k}$; $M$, $T_{\max}$, $\delta$, $\sigma^2_{\min}$, seed
\State $\pi_m\gets 1/M$; $\mu_m\gets$ distinct samples of $Z_k$; $\sigma^2_m\gets\max(\mathrm{var}(Z_k),\sigma^2_{\min})$
\State $\mathcal{L}^{(0)}\gets-\infty$
\For{$t=1,\dots,T_{\max}$}
  \State E-step: $\log r_{im}\gets\log\pi_m+\log\mathcal N(z_i;\mu_m,\mathrm{diag}\,\sigma^2_m)-\operatorname{logsumexp}_{m'}[\cdot]$
  \State M-step: $N_m\gets\sum_i r_{im}$; update $\pi_m,\mu_m,\sigma^2_m$ (variance floored at $\sigma^2_{\min}$)
  \State $\mathcal{L}^{(t)}\gets\frac{1}{N_k}\sum_i\log p(z_i)$; \textbf{if} $|\mathcal{L}^{(t)}-\mathcal{L}^{(t-1)}|<\delta$ \textbf{break}
\EndFor
\State Final E-step to obtain $r_{im}$
\State $S_k\gets\sum_i\sum_m r_{im}(z_i-\mu_m)(z_i-\mu_m)\T$ \Comment{float64}
\State \Return $(\pi,\mu,\sigma^2)$, $S_k$, $N_k$
\end{algorithmic}
\end{algorithm}

\begin{algorithm}[h]
\caption{\textsc{LearnTask}$(k, D_k)$}\label{alg:learn}
\begin{algorithmic}[1]
\Require persistent state (Section~\ref{sec:state}); training data $D_k$ of the current task only
\State Split $D_k$ into $D_k^{\mathrm{tr}}$ and held-out $H_k$ (5\%, $\le 500$, fixed seed)
\State $Z_k\gets\{z(x): x\in\text{sample}_{\le N_{\max}}(D_k^{\mathrm{tr}})\}$ \Comment{adapters disabled}
\If{$k\ge2$}
  \State $a_k\gets$ affinity of $Z_k$ to tasks $1..k-1$ (router state \emph{before} registering $k$)
  \State $\beta^{(\ell)}_{k,1:k-1}\gets\sum_{i<k}a_{ki}\,C^{(\ell)}_{i,1:k-1}$ for all $\ell\in\Ladapt$
\EndIf
\State $(\pi_k,\mu_k,\sigma^2_k),S_k,N_k\gets$ \Call{FitTaskGMM}{$Z_k$}
\State Register GMM $k$; $S\gets S+S_k$; $n\gets n+N_k$; recompute all Cholesky factors
\If{$k\ge2$} calibrate $\tau_k$ on $H_k$ and update $\tau$ \EndIf
\State Add atom $k$ ($A_k$ Kaiming-uniform, $B_k=0$); freeze atoms $1..k-1$
\If{$k\ge2$} \textbf{Phase A}: train $\beta_k$ only on $\mathcal L_k$ for $S_A$ steps ($B_k$ held at 0) \EndIf
\State \textbf{Phase B}: train $(A_k,B_k)$ and $\beta_k$ on $\mathcal L_k$ for 3 epochs
\State Freeze atom $k$, row $k$ of all $C^{(\ell)}$, GMM $k$; write checkpoint and hashes
\State Discard $D_k$, $Z_k$, $H_k$ \Comment{Invariant I4}
\end{algorithmic}
\end{algorithm}

\begin{algorithm}[h]
\caption{\textsc{Predict}$(x, n)$ --- the only inference entry point}\label{alg:predict}
\begin{algorithmic}[1]
\Require prompt string $x$; number of samples $n$ \Comment{no task metadata (I1)}
\State $z\gets z(x)$; compute $s_k(z)$ for $k=1..K$ (Eq.~\ref{eq:score})
\State $\alpha\gets e_{\hat k}$ (hard) or $\mathrm{softmax}(s/\tau)$ (soft)
\State $c^{(\ell)}\gets C^{(\ell)\T}\alpha$ for all $\ell\in\Ladapt$; set atom-bank coefficients
\State Generate $n$ samples in one round \Comment{no execution, no second round (I2, I3)}
\State \Return samples, $\alpha$, $\hat k$
\end{algorithmic}
\end{algorithm}

\begin{algorithm}[h]
\caption{\textsc{RunSequence}$(T_1,\dots,T_K)$}\label{alg:run}
\begin{algorithmic}[1]
\State Initialise backbone (frozen), $P$ (seed $s_P$), $S\gets0$, $n\gets0$, empty registries
\For{$k=1,\dots,K$}
  \State \Call{LearnTask}{$k, D_k$}
  \LineComment{Evaluation only: read-only access to test sets of tasks $1..k$}
  \For{$j=1,\dots,k$}
     \State generate with \Call{Predict}{} for all test prompts of $T_j$; write to disk
     \State score with the evaluation harness $\rightarrow a_{k,j}$; record hard-routing confusion row
  \EndFor
\EndFor
\end{algorithmic}
\end{algorithm}

\clearpage
% =====================================================================
\section{Hyperparameters}\label{sec:hparams}
% =====================================================================

\begin{table}[h]
\centering
\small
\begin{tabularx}{\textwidth}{@{}l X l l@{}}
\toprule
Symbol & Meaning & Default & Source\\
\midrule
--- & Backbone & Qwen2.5-Coder-1.5B & thesis\\
$\ell^*$ & Routing/split block & 4 & thesis\\
$p$ & Routing dimension & 256 & thesis\\
$s_P$ & Projection seed & 0 & new\\
--- & Max prompt length (routing), truncation side & 2048, right & thesis\\
$M$ & GMM components per task & 4 & thesis\\
$T_{\max}$, $\delta$ & EM iterations, tolerance & 50, $10^{-4}$ & thesis\\
$\sigma^2_{\min}$ & Variance floor & 0.02 & thesis\\
$N_{\max}$ & Embeddings per task for GMM & 5000 & thesis\\
$\lambda$ & Covariance shrinkage & 1.0 (sweep 0, .25, .5, .75, 1) & new\\
$\epsilon_\Sigma$ & Covariance jitter & $10^{-4}$ & new\\
$\tau_{\mathrm{aff}}$ & Temperature for affinity & 1.0 & new\\
$\tau$ & Routing temperature & calibrated on $H_k$, grid $2^{-2}..2^{8}$ & new\\
$|H_k|$ & Held-out slice for calibration & 5\% of $D_k$, $\le 500$ & new\\
$r$, \texttt{lora\_alpha}, dropout & LoRA atom & 16, 32, 0.1 & thesis\\
--- & Target modules, blocks & \texttt{q\_proj}, \texttt{v\_proj}; $\ell^*..L-1$ & thesis\\
$S_A$, $\eta_A$ & Phase A steps, LR ($\beta$ only) & 300, $5\times10^{-3}$ & new\\
--- & Phase B epochs, batch, LR, warmup, WD (atom) & 3, 32, $10^{-4}$, 0.01, 0.01 & thesis\\
$\eta_\beta$ & Phase B LR for $\beta$ & $2\times10^{-5}$ & new\\
--- & Max train / val / test examples per task & 100k / 2k / 5k & thesis\\
--- & CodeTask decoding & $T{=}0.2$, top-$p$ 0.95, 512 tokens, hard routing & thesis\\
--- & Executable decoding & $n{=}5$, 2048 tokens; $T{=}0.2$, top-$p$ 0.95 & thesis / assumed\\
--- & Precision & bf16 backbone, fp32 atoms/$\beta$, fp64 router & new\\
\bottomrule
\end{tabularx}
\caption{Hyperparameters. Every run \MUST\ save the full resolved configuration with its checkpoints.}
\end{table}

% =====================================================================
\section{Software architecture}\label{sec:software}
% =====================================================================

The following layout and signatures are recommended (\SHOULD). Names may change; responsibilities and the
inference signature may not.

\begin{lstlisting}
cgr/
  config.py        # CGRConfig dataclass: every value of Section 7; saved as config.json
  data.py          # task registry; reuse thesis loaders/templates; split D_k -> (train, H_k)
  features.py      # RoutingEncoder: P, layer l*, masked mean pool, parameter-free LN
  gmm.py           # fit_diag_gmm (Algorithm 1), within_scatter
  router.py        # Router: register, scores, weights, affinity, calibrate_tau
  atoms.py         # AtomBankLinear, install_atom_bank, set_coefficients
  composition.py   # Composition: per-block lower-triangular C, init_row, freeze_row
  train.py         # learn_task (Algorithm 2): phase A, phase B
  infer.py         # predict (Algorithm 3) -- the ONLY inference entry point
  run_sequence.py  # Algorithm 4: loop over tasks, checkpoints, evaluation schedule
  eval/            # reuse thesis BLEU/CodeBLEU/Pass@k harness; cl_metrics.py
  tests/           # acceptance tests T1-T10 (Section 10)
\end{lstlisting}

\begin{lstlisting}
class RoutingEncoder:
    def embed(self, prompts: list[str]) -> np.ndarray: ...       # [B, p], adapters disabled

def fit_diag_gmm(Z: np.ndarray, M: int, t_max: int, tol: float,
                 var_floor: float, seed: int) -> tuple["DiagGMM", np.ndarray]: ...  # (gmm, resp)
def within_scatter(Z: np.ndarray, gmm: "DiagGMM", resp: np.ndarray) -> np.ndarray: ...  # [p, p]

class Router:
    def register(self, gmm: "DiagGMM", scatter: np.ndarray, n_k: int) -> None: ...
    def scores(self, Z: np.ndarray) -> np.ndarray: ...            # [B, K], Eq. (score)
    def weights(self, Z: np.ndarray, hard: bool) -> np.ndarray: ...  # [B, K]
    def affinity(self, Z_new: np.ndarray) -> np.ndarray: ...      # [K_prev], before register()
    def calibrate_tau(self, Z_heldout: np.ndarray, k: int) -> float: ...

class AtomBank:
    def add_atom(self) -> int: ...
    def freeze_atom(self, i: int) -> None: ...
    def set_coefficients(self, c: torch.Tensor) -> None: ...      # [B, n_blocks, K]

class Composition:
    def init_row(self, k: int, affinity: np.ndarray) -> None: ...
    def trainable_row(self, k: int) -> torch.nn.Parameter: ...    # [n_blocks, k-1]
    def freeze_row(self, k: int) -> None: ...
    def coefficients(self, alpha: torch.Tensor) -> torch.Tensor: ...  # C^T alpha -> [B, n_blocks, K]

def learn_task(k: int, task_data, state, cfg) -> None: ...
def predict(prompts: list[str], n_samples: int, state, cfg) -> list["Prediction"]: ...
    # MUST NOT accept task ids, languages, tests or any other metadata
\end{lstlisting}

\paragraph{Checkpoint layout.} After each task:
\begin{lstlisting}
state/
  config.json
  projection_P.npy
  router/gmm_{k}.npz          # pi [M], mu [M,p], var [M,p]
  router/scatter.npz          # S [p,p], n, S_k for every k
  router/tau.json             # tau_k for every k and the running tau
  atoms/atom_{k}.safetensors  # A_k, B_k for every (block, module)
  composition/beta_{k}.npy    # [n_blocks, k-1]
  manifest.json               # sha256 of every frozen artifact (test T7)
\end{lstlisting}

% =====================================================================
\section{Evaluation protocol}\label{sec:eval}
% =====================================================================

\paragraph{Task sequences\tagthesis.} CodeTask: CONCODE, CodeTrans, CodeSearchNet, BFP, KodCode, RunBugRun,
TheVault, CoST. Executable: Python, C++, C\#, Java, Rust, Swift, PHP, TypeScript, Shell.

\paragraph{Per-task metrics\tagthesis.} CodeTask: BLEU for text-output tasks (CodeSearchNet, TheVault) and
CodeBLEU for code-output tasks, as in the thesis evaluation code\tagassumed. Executable: Pass@1 and Pass@5 with the
unbiased estimator $\mathrm{Pass}@k=1-\binom{n-c}{k}/\binom{n}{k}$, $n=5$.

\paragraph{Continual-learning metrics\tagthesis.} With $a_{i,j}$ the score on task $j$ after training tasks
$1..i$:
\begin{gather}
  \mathrm{FM}_K=\frac{1}{K-1}\sum_{j=1}^{K-1}\max_{i<K}(a_{i,j}-a_{K,j}),\qquad
  \mathrm{BWT}_K=\frac{1}{K-1}\sum_{j=1}^{K-1}(a_{K,j}-a_{j,j}),\\
  \mathrm{FWT}_K=\frac{1}{K-1}\sum_{j=2}^{K}(a_{j,j}-\tilde a_j).
\end{gather}
$\tilde a_j$ \MUST\ match the reference used in the thesis code\tagassumed. In addition, report the transfer gain
attributable to composition, $\Delta_j=a_{j,j}^{\mathrm{CGR}}-a_{j,j}^{C=I}$, and the Phase-A-only score of each
task (pure reuse).

\paragraph{Routing diagnostics.} Row-normalised hard-routing confusion matrices after the final task and after
each task; per-task routing accuracy over the sequence (to measure routing drift caused by updates of $\bar\Sigma$);
mean entropy of $\alpha(x)$; the learned $\beta$ matrices averaged over blocks.

\paragraph{Reference numbers (thesis, Qwen2.5-Coder-1.5B).} The leak-free GMM-Router baseline on the executable
benchmark is \emph{Round-1 only}.
\begin{center}
\small
\begin{tabular}{@{}lccc@{}}
\toprule
Method & CodeTask avg (\%) & Exec Pass@1 (\%) & Exec Pass@5 (\%)\\
\midrule
GMM-Router, leak-free (Round-1 only) & 51.62 & 33.04 & 41.66\\
GMM-Router + Round-2 (leaky, not a baseline) & --- & 34.40 & 47.75\\
LoRA per task (oracle task identity) & 54.15 & 33.46 & 42.19\\
SeqLoRA & 34.90 & 37.51 & 45.67\\
TreeLoRA & 34.58 & 36.84 & 46.37\\
O-LoRA & 42.77 & 33.71 & 42.09\\
\bottomrule
\end{tabular}
\end{center}
Hard-routing accuracy of the thesis router: 90.6\% (CodeTask), 15.8\% (executable).

\subsection{Ablation grid}\label{sec:ablations}
\begin{center}
\small
\begin{tabular}{@{}llll@{}}
\toprule
ID & Composition & $\lambda$ & Purpose\\
\midrule
A0 & $C=I$ & 0 & Reproduce thesis Round-1 (sanity check)\\
A1 & $C=I$ & 1 & Routing only\\
A2 & learned & 0 & Transfer only\\
A3 & learned & 1 & Full CGR\\
A4 & learned & 0.25, 0.5, 0.75 & Shrinkage sweep\\
A5 & Phase A only (atom $k$ never trained) & 1 & Pure reuse of earlier knowledge\\
A6 & learned, no Phase A & 1 & Value of compose-first training\\
A7 & learned, $\beta$ initialised at 0 & 1 & Value of the router prior\\
A8 & learned, oracle routing $\alpha=e_{\text{true}}$ & --- & Composition upper bound (analysis only)\\
A9 & full CGR, fixed $\tau\in\mathcal T$ & 1 & Temperature sensitivity (analysis only)\\
\bottomrule
\end{tabular}
\end{center}
A8 and A9 use test-set information and \MUSTNOT\ be presented as methods. Report at least 3 seeds for A0--A3;
on the executable benchmark one problem is 2 points of a language's Pass@1 (50 problems per language), so report
bootstrap confidence intervals over problems.

% =====================================================================
\section{Acceptance tests}\label{sec:tests}
% =====================================================================

\begin{enumerate}[label=T\arabic*., leftmargin=*]
  \item \textbf{Adapter-independent features.} With random non-zero atoms and coefficients installed, $z(x)$ read
        from the full model equals $z(x)$ computed with adapters disabled (max abs diff $<10^{-4}$, fp32).
  \item \textbf{EM reproduction.} On synthetic data from a known diagonal GMM, the fitted means are recovered
        (up to permutation) and the log-likelihood is non-decreasing across iterations.
  \item \textbf{Scatter correctness.} With $M=1$ on synthetic labelled data, $S/n$ equals the directly computed
        pooled within-class covariance (relative tolerance $10^{-6}$).
  \item \textbf{Thesis router equivalence.} With $\lambda=0$, \texttt{Router.scores} equals the diagonal-GMM
        log-likelihood plus $\log(1/K)$ computed independently (thesis Eq.~4.6), to $10^{-8}$.
  \item \textbf{Merge equivalence.} For random $\alpha$ and lower-triangular $C$, the atom-bank output with
        $c=C\T\alpha$ equals $Wx+\kappa\sum_k\alpha_k\Delta W_k x$ computed from explicitly merged matrices
        (relative tolerance $10^{-4}$). With $C=I$ this equals thesis Eq.~4.10.
  \item \textbf{Phase continuity.} In eval mode (dropout off), the held-out loss at Phase~B step~0 equals the
        Phase~A final held-out loss (difference $<10^{-5}$).
  \item \textbf{Freeze integrity.} After \textsc{LearnTask}$(k)$, the sha256 hashes of atoms $1..k-1$, rows
        $2..k-1$ of every $C^{(\ell)}$, GMMs $1..k-1$ and $P$ are unchanged.
  \item \textbf{Leakage guard.} \texttt{predict} has no parameter other than prompts and $n$; a static check
        confirms that no module under \texttt{cgr/} other than \texttt{eval/} imports \texttt{subprocess},
        the execution harness or test-case loaders; the evaluation harness runs only on files written after
        generation has finished.
  \item \textbf{No-replay guard.} The data registry refuses to construct a loader for task $j<k$ while training
        task $k$; only evaluation code may open earlier test splits, in read-only mode.
  \item \textbf{Determinism.} With fixed seeds and greedy decoding, two runs produce identical $\alpha$, $\beta$
        and generations.
\end{enumerate}

% =====================================================================
\section{Known risks and points to verify}\label{sec:risks}
% =====================================================================

\begin{itemize}[leftmargin=*]
  \item \textbf{Routing drift from $\bar\Sigma$ updates.} Registering a task changes $\bar\Sigma$ and therefore
        earlier tasks' scores. Monitor per-task routing accuracy after every task; a variant that freezes
        $\bar\Sigma$ after the first three tasks \MAY\ be added if drift is large.
  \item \textbf{Input similarity is not output compatibility.} The affinity prior can point to unhelpful atoms;
        Phase~A is responsible for correcting it. Compare Phase-A loss with the loss of a fresh LoRA after the same
        number of steps to detect negative transfer.
  \item \textbf{First task.} Task 1 receives no forward transfer by construction.
  \item \textbf{Temperature scale.} If calibrated $\tau$ values differ by more than an order of magnitude across
        tasks, report them and the A9 sensitivity.
  \item \textbf{Points to verify against the thesis code} (all tagged \tagassumed): hidden-state index convention
        for $\ell^*$; executable sampling temperature and top-$p$; the reference $\tilde a_j$ used for FWT; the
        exact BLEU/CodeBLEU assignment per CodeTask task.
\end{itemize}

\paragraph{Related work to position against.} Progressive Neural Networks (lateral connections to frozen
columns), LoraHub (composition of frozen LoRA modules) and MoCL (rehearsal-free composition of task modules for
language models). CGR differs in that task identification (GMM routing over tasks) and knowledge composition
(the learned matrix $C$) are learned separately and meet only in $c(x)=C\T\alpha(x)$.

\end{document}